% =======> To use this template:
% => Edit authors in this file.
% => Edit the paper's content itself in the respective sections. Edit as needed
% => Add any references to references.bib (cite with \cite or \citet)
% =================================================

\documentclass{article}

% to avoid loading the natbib package, add option nonatbib:
% \usepackage[nonatbib]{neurips_2024}
% \PassOptionsToPackage{square,comma,sort,numbers}{natbib}
\PassOptionsToPackage{numbers, compress}{natbib}
% \usepackage[numbers]{natbib}

% to compile a preprint version, e.g., for submission to arXiv, add add the
% [preprint] option:
\usepackage[preprint]{neurips_2024}
% \usepackage[nonatbib, preprint]{neurips_2024}

\usepackage[utf8]{inputenc} % allow utf-8 input
% \usepackage[T1]{fontenc}    % use 8-bit T1 fonts
% For maths/theorems and such
\usepackage{amsmath}
\usepackage{amssymb}
\usepackage{mathtools}
\usepackage{amsthm}
\usepackage{float}
% Use DeclareMathOperator and/or \newcommand to define abbreviations, e.g.:
\DeclareMathOperator{\R}{\mathbb{R}} % real numbers
% Recommended, but optional, packages for figures and better typesetting:
\usepackage{microtype}
\usepackage{graphicx}
\usepackage{subfigure}
\usepackage{subcaption}
\usepackage{booktabs} % for professional tables (https://www.ctan.org/pkg/booktabs)
\usepackage{algorithm}  % For algorithms
\usepackage{algorithmic}
\usepackage{hyperref} % hyperref makes hyperlinks in the resulting PDF.
\usepackage{wrapfig}  % For wrapping figures with text
\usepackage[capitalize,noabbrev]{cleveref}   % For \cref
\RequirePackage{doi}
% Color references and URLs
\usepackage{xcolor, colortbl}
\definecolor{darkblue}{rgb}{0.19, 0.19, 0.62}
\hypersetup{
  colorlinks = true,
  citecolor  = darkblue,
  filecolor  = darkblue,
  urlcolor   = darkblue,
    % linkcolor  = myred,
}
% Todonotes is useful during development; simply uncomment the next line
%    and comment out the line below the next line to turn off comments
%\usepackage[disable,textsize=tiny]{todonotes}
\usepackage[textsize=tiny]{todonotes}

% \input{_macros}



\title{
CSE 151B Milestone Report
\\
{\small \url{https://github.com/SatrunsDream/math-qa-llm}} % Note: You can put it in your abstracts or as a footnote in the first page, if you prefer. Note 2: Use \href or \url to link your URL.
}

\author{
Tia Irani\\
A17303637\\
\And
Leonel Alejandro Montoya\\
A17525034 \\
\And
Diego Osborn \\
A17676487 \\
\And
Sardor Sobirov\\
A17656802\\
\and
Team: Probably Overfitting
}
\begin{document}
\maketitle
\begin{abstract}
% Check out \href{https://www.alignmentforum.org/posts/eJGptPbbFPZGLpjsp/highly-opinionated-advice-on-how-to-write-ml-papers#Abstract}{this blog post for tips on how to write a compelling abstract and paper.} \textbf{\textit{Include your final private Kaggle leaderboard score and/or rank at the end of the abstract or introduction.}}
We develop an inference-time optimization pipeline for mathematical reasoning using the pretrained Qwen/Qwen3-4B-Thinking-2507 model. Competition constraints fix the model weights and prohibit external tools, so we focus on structured prompting, adaptive multi-phase inference, and evaluation-aware output formatting. On a 50-question public subset, our adaptive system achieves 42\% overall accuracy (55\% MCQ, 33.3\% free-form), up from 38\% with prompting alone. These results are preliminary---reflecting both the task's difficulty for a 4B model and limited evaluation scale---but the 15 percentage point MCQ gain demonstrates viable inference-time improvement without weight modification.
\end{abstract}

% \section*{Formatting instructions (remove this section in your submission)}
% Please read the instructions below carefully and follow them faithfully. You can keep the section headers but delete the instructional texts.
% \paragraph{Hints on writing}
% \begin{itemize}
%     \item Each paragraph should convey one main message, typically in the beginning sentence of the paragraph (topical sentence).
%     \item  Use simple sentences with fewer compounds or clauses. The simpler, the better. Avoid passive tense and use active tense.
%     \item Use PDF for your figures to ensure the resolution of your images.
%     \item The captions of Figures and Tables should contain sufficient details to be self-explanatory (readers can understand the figure without referring to the text).
% \end{itemize}


% Figures and tables must be appropriately referenced and discussed in the main text (tip: Use \textbackslash label\{fig:xyz\} and \textbackslash cref\{fig:xyz\}; example:~\cref{fig:example}). Captions must strive to be self-contained.
% \paragraph{Figures.}
% % \begin{figure}[h]  % Use figure or wrapfigure as needed.
% \begin{wrapfigure}{r}{0.4\textwidth} % Example: occupy 40% of text width
%   \vspace{-4mm}  % Force moving the figure up by a little bit
%   \centering
%   \fbox{\rule[-.5cm]{0cm}{4cm} \rule[-.5cm]{4cm}{0cm}}
%   \caption{Sample figure caption.}
%   \label{fig:example}
% \end{wrapfigure}
% % \end{figure}


% All artwork must be neat, clean, and legible. Lines should be dark enough for
% purposes of reproduction. The figure number and caption always appear after the figure. The figure caption should be lowercase (except for the first word and
% proper nouns); figures are numbered consecutively. Captions should be written clearly and, ideally, allow the reader to understand the key message of the figure without having to read the main text.


% \paragraph{Tables.}
% All tables must be centered, neat, clean, and legible.  The table number and
% title always appear before the table. See~\cref{sample-table} for an example.
% Captions should be written clearly and, ideally, allow the reader to understand the key message of the table without having to read the main text.

% % % === Introduction
\section{Introduction}
\label{sec:intro}

\paragraph{Problem Definition.}

This project addresses mathematical reasoning using a large language model (LLM). Given a set of math problems, the LLM must generate adequate answers to each. The mathematical questions include multiple-choice and free-form problems, where the model must compute a correct answer.

\paragraph{Problem Significance.}

Mathematical reasoning is an important benchmark for deep learning because it requires understanding complex questions, being able to perform multi-step logic, computing numerical or symbolic answers, and presenting answers in a structured format. Improving these abilities can help build AI systems that are more capable of aiding humans in engineering, scientific research, and education. 

A reliable model capable of mathematical reasoning could assist engineers in verifying calculations that are safety-critical, process large amounts of data for scientists, such as removing noisy measurements, and even explain steps in mathematical calculations to students. 

\paragraph{Technical Challenge.}

The task is difficult because mathematical problems often require multiple steps of computation and logic, and many times don't have a single correct procedure. The model must handle problems being presented in multiple formats, reason through the multiple steps in computation that may be involved, and provide the best answer when there may be multiple correct choices. Overall, this brings multiple challenges to the processing of data for the model to understand, designing a model that generalizes well in the broad field of mathematics, and then having the model present answers in a format readable to humans. 

Secondly, there still exists the challenge of providing enough computing resources to the model for computation and evaluation. To find a reliable model capable of mathematical reasoning, it must be tested both extensively and intensively to get a full understanding and confidence in the model's capabilities
. 

\paragraph{Contributions.}

% Find and cite (use \href{https://www.overleaf.com/learn/latex/Bibliography_management_with_bibtex}{BibTex
% }) state-of-the-art works or papers that have tried to solve a similar problem (use \href{https://scholar.google.com/}{Google Scholar} to search). Point out the limitations of the state-of-the-art.

Our main contributions are:
\begin{itemize}
  \item \textbf{Contribution 1: Baseline characterization and failure analysis.}
We first establish a baseline using the unmodified pipeline on a random subset of 10 public questions. This initial run achieved:
\begin{itemize}
    \item MCQ accuracy: 50\% (2/4)
    \item Free-form accuracy: 33.33\% (2/6)
    \item Overall accuracy: 40\% (4/10)
\end{itemize}

Analysis of individual failures revealed three dominant error modes:
\begin{itemize}
    \item Truncation of long reasoning traces, leading to missing final answers.
    \item Incorrect formatting for multi-part free-form answers, causing zero credit despite partially correct reasoning.
    \item Numerical precision mismatches (e.g., correct value but slightly different representation).
\end{itemize}

These observations guided all subsequent design decisions.

\item \textbf{Contribution 2: Structured prompting}
We replace the baseline prompt with a structured reasoning template:
\textit{UNDERSTAND $\rightarrow$ PLAN $\rightarrow$ SOLVE $\rightarrow$ VERIFY $\rightarrow$ ANSWER}.
This explicitly encourages step-by-step decomposition of problems, verification of intermediate results, and consistent final answer formatting.

Additionally, we introduce explicit instructions for:
\begin{itemize}
    \item Ordering multiple answers for \texttt{[ANS]} placeholders
    \item Formatting MCQ outputs with a final answer of the form \texttt{\textbackslash boxed\{X\}}
\end{itemize}

This reduces evaluation failures caused by ambiguous or unparsable outputs and aligns model behavior with scoring rules.

\item \textbf{Contribution 3: Adaptive inference to balance accuracy and runtime.}
We introduce an adaptive inference pipeline \cite{snell2024scaling} that allocates compute selectively:
\begin{itemize}
    \item Phase 1: low-cost generation (\texttt{max\_tokens $\approx$ 2048}, single sample) for all questions
    \item Phase 2: higher-budget retry (\texttt{max\_tokens $\approx$ 6144}, multi-sample) for uncertain outputs
    \item Phase 3: full-budget reasoning (\texttt{max\_tokens $\approx$ 8192}, majority voting) for remaining difficult cases
\end{itemize}

Uncertainty is detected using signals such as:
\begin{itemize}
    \item Truncated outputs (\texttt{finish\_reason = length})
    \item Missing \texttt{\textbackslash boxed\{\}} answers
    \item Weak agreement across multiple samples
\end{itemize}

This approach avoids the high cost of applying full self-consistency \cite{wang2022selfconsistency} to all problems while improving reliability on harder instances.
\item \textbf{Contribution 4: Handling long reasoning traces in thinking models.}
We empirically observed that Qwen3-Thinking \cite{qwen3} often produces thousands of tokens before reaching a final answer. Using a smaller generation budget, such as \texttt{max\_tokens = 2048}, can truncate the response before the model produces a final \texttt{\textbackslash boxed\{\}} answer.

To reduce this failure mode, we increased the generation budget to \texttt{max\_tokens = 8192}. This improves answer extractability while balancing runtime cost.

\item \textbf{Contribution 5: Evaluation-aware output design and dataset-specific handling.}
We adapted the pipeline to the structure of the competition dataset. The system separately handles multiple-choice and free-form questions, preserves the full model response for submission, and encourages final answers to appear in a parseable \texttt{\textbackslash boxed\{\}} format.

For free-form questions, the pipeline also accounts for cases with multiple \texttt{[ANS]} placeholders, where all sub-answers must be correct for the item to count as correct. For multiple-choice questions, the prompt asks the model to place the final selected option in the form \texttt{\textbackslash boxed\{X\}}.

\item \textbf{Contribution 6: System-level optimization for scalable inference.}
We improved scalability by moving from a slower sequential Transformers-based workflow to a vLLM-based \cite{kwon2023vllm} backend under Windows Subsystem for Linux 2 (WSL 2) with GPU support. In earlier runs, sequential generation could take several minutes per question, which would make full public or private inference impractical.

Using vLLM improves GPU utilization and supports more efficient long-context generation, making large-scale inference on the full dataset more feasible.

\item \textbf{Contribution 7: Reproducible experimentation and validation workflow.}
We extended the starter notebook into a more reproducible experiment pipeline. The workflow supports random public validation batches, JSONL logging of predictions and correctness, checkpointing for longer runs, and automatic CSV generation for private leaderboard submission.

This makes it easier to compare prompt and inference changes across runs, avoid stale notebook outputs, and reliably transition from public validation to private inference.
    % etc.
\end{itemize}

The starter code provides the baseline workflow for the competition. It starts with loading the public dataset, separates multiple-choice and free-form questions, constructs prompts for each type, loads the model and runs inference with vLLM, scores the responses and reports accuracy, and finally saves the results in JSON format for submission.  

% % % % === Related Work
% \section{Related Work (Optional)}
% \label{sec:relatedwork}
% Describe related papers and contextualize your own work with respect to them. If useful, you can reference and discuss them in the introduction too. Example:
% \paragraph{Related works topic 1.} \citet{radford2018improving} introduced the first GPT-1 model...  % This is an example for citing (plz remove)

% \paragraph{Related works topic 2.}
% % 

% % % === Methods
\section{Methods}
\label{sec:methods}
In this section, clearly describe the technical details of the problem and your approach. If needed, provide the necessary background and cite any sources that you used.


\paragraph{Problem Setting.}
% {\textcolor{red}{Introduce the problem. Formulate your prediction task as an optimization problem and write out the objective (i.e., loss function) that you use to train your model. Describe the dataset. Define the inputs and outputs in mathematical language, describing their meaning in natural language. \textbf{Note:} \textit{The inputs and outputs here should correspond to what your model actually ingests and predicts, respectively}.}


This project studies mathematical reasoning using a large language model (LLM). Given a dataset of math problems, the goal is to generate correct answers for both multiple-choice (MCQ) and free-form questions. The problems span a wide range of difficulty (high-school to graduate level mathematics) and are written in LaTeX, thus requiring both symbolic understanding as well as multi-step reasoning.


% {\textcolor{blue}{\textbf{Update tomorrow with more information:}} 
We formulate this as a supervised learning problem. Let the dataset be:
\[\mathcal{D} = \{(x_i, y_i)\}_{i=1}^N\]

where:
\begin{itemize}
    \item $x_i$ is a math problem written in natural language
    \item $y_i$ is the corresponding correct answer
\end{itemize}

The model, parametrized by $\theta$,  defines a conditional distribution:
\[P_\theta(y|x)\]
Training conceptually minimizes the negative log-likelihood and encourages the model to assign high probability to correct answers:
\[\mathcal{L}(\theta) = - \sum_{i=1}^{N}\log P_\theta(y_i | x_i)\]
For this project, instead of training from scratch, we use a pretrained reasoning model, Qwen/Qwen3-4B-Thinking-2507 \cite{qwen3}. Optimization of performance is done via prompting and inference strategies rather than gradient-based updates.

\paragraph{Datasets}
The dataset is provided in .jsonl formatting, where each line is a JSON object representing a single problem. The only available files are:
\begin{itemize}
    \item \textbf{public.jsonl}: labeled, and used for validation.
    \item \textbf{private.jsonl}: unlabeled, and used for submission
\end{itemize}
Each JSON object includes the following four fields to represent a problem:
\begin{itemize}
    \item \textbf{id:} Unique integer identifier for the problem.
    \item \textbf{question:} Problem statement written in LaTeX. Free-form items use [ANS] in place of where answers belong.
    \item \textbf{answer:} Ground truth for the problem. Free-form responses will have a list of strings, one string per [ANS] present in the question field. Multiple-Choice will have a single capital letter for the correct answer.
    \item \textbf{options:} A list of candidate choices in LaTeX. Options will only be present for multiple-choice questions. 
\end{itemize}

There are two task types:
\begin{enumerate}
    \item \textbf{Free-form:} Free-form questions will require the model to produce one or more numerical or symbolic answers matching each [ANS]. 
    % \begin{quote}
    %     \textit{Example Input:}
    %     \\\{``question": ``Here is an expression with negative exponents.\textbackslash n\$\textbackslash\textbackslash frac\{1\}\{(-8)\^\{-3\}\}=\$ [ANS]\textbackslash n Evaluate the expression.",``answer": ["-512"],``id": 4\}
    % \end{quote}
    \item \textbf{Multiple-Choice:} Multiple-choice questions will require the model to select the correct option. In order for the item to be marked correct, the correct letter must be chosen.
    % \begin{quote}
    %     \textit{Example Input (abbreviated):}
    %     \\\{``question": ``Given \$u(x, y) = x\^3 + ...\$ find \$f(z)\$ ...", ``options": ["\$\$...\$\$", ``\$\$...\$\$"], ``answer": ``C", ``id": 1\}
    % \end{quote}
\end{enumerate}

A key property of the dataset is that free-form questions require exact matching of all sub-answers. This requirement makes them stricter than standard generation tasks.

\paragraph{Inputs and Outputs} 
Inputs $x \in \mathcal{X}$ are sequences of tokens representing a problem statement in the natural language. Outputs $y \in \mathcal{Y}$ are either discrete labels or token sequences, for multiple-choice and free-form questions respectively). In natural language, the input is the problem statement and the output is the final answer(s).

\paragraph{Idea Summary.}
% \textcolor{red}{Describe and motivate your idea to solve the problem.  How is it different from the state-of-the-art method? Provide a sketch of your idea to illustrate y our idea. Is there any risk associated with your idea? Provide alternative solutions if your idea does not work out.}

Instead of modifying model weights, our approach improves reasoning performance through inference-time techniques tailored to the constraints of the competition.

\subparagraph{Structured Prompting} Due to the baseline prompt producing formatting errors and incomplete reasoning traces, we replaced a simple prompt with a structured reasoning template:
\begin{center}
    \textbf{UNDERSTAND} $\rightarrow$ \textbf{PLAN} $\rightarrow$ \textbf{SOLVE} $\rightarrow$ \textbf{VERIFY} $\rightarrow$ \textbf{ANSWER}
\end{center}

This structure encourages the model to:
\begin{enumerate}
    \item Parse the problem carefully,
    \item Plan a solution strategy,
    \item Execute step-by-step reasoning,
    \item Verify intermediate results,
    \item Produce a clearly formatted final answer.
\end{enumerate}

This follows the chain-of-thought paradigm \citet{wei2022cot}, which demonstrates that explicit step-by-step reasoning improves LLM performance on mathematical tasks. As a result, common failure modes such as incomplete reasoning, arithmetic errors, and missing final outputs are reduced.

\subparagraph{Adaptive Multi-Pass Inference} Additionally, rather than using fixed compute for every question, we use a multi-phase strategy:
\begin{enumerate}
    \item \textbf{Phase 1:} fast, low-token reasoning
    \item \textbf{Phase 2:} deeper reasoning for uncertain cases
    \item \textbf{Phase 3:} high-budget reasoning and multiple samples
\end{enumerate}

Uncertainty is detected by heuristics such as:
\begin{itemize}
    \item the absence of a final answer (by missing \textbackslash boxed\{\}),
    \item truncated outputs,
    \item and inconsistent or low-confidence responses.
\end{itemize}
This approach balances accuracy and efficiency by reserving expensive computation for difficult problems.

\subparagraph{Self-Consistency} Following \cite{wang2022selfconsistency}, we optionally generate multiple outputs:
\[\{y^{(1)},y^{(2)},y^{(3)},\dots,y^{(K)}\}\]
Out of these outputs, we select the final prediction $\hat{y}$ via majority voting. While this method often improves robustness, it significantly increases runtime and is therefore disabled for large-scale runs. 

\subparagraph{Output Formatting} To ensure compatibility with the evaluation pipeline, we explicitly instruct the model to output final answers in standardized format (e.g., \textbackslash boxed\{\} and sorted for free-form questions). Because the evaluation pipeline extracts answers from the raw response string and not a separate prediction field, this output formatting is a critical constraint.

\subparagraph{Limitations and Fallback Strategies} Despite the effectiveness of structured prompt engineering and adaptive inference, our approach has several limitations stemming from both the dataset and the constraints of the competition setting.
\begin{itemize}
    \item \textbf{High Computational Cost:} The use of a "thinking" model introduces significant computational overhead. Thousands of tokens are often required per question in order to generate long reasoning traces, which in turn means that, per question, each inference takes minutes to be generated. Running the model on a dataset like public.jsonl, which has 1126 problems, can take several hours to days. This is especially true when multi-pass inference or self-consistency is enabled. As such, the extent of hyperparameter tuning and large-scale experimentation is limited.
    \item \textbf{Sensitivity to Prompt Design:} Because the model parameters are fixed, performance depends heavily on prompt engineering. Small changes in wording, structure, or formatting instructions can lead to large variations in output quality. The result of such sensitivity is a brittle system that is difficult to systematically optimize. This is primarily due to the fact that improvements are not guaranteed to generalize across different problem types.
    \item \textbf{Output Formatting Dependency:} As stated previously the evaluation pipeline extracts answers directly from the model's raw output. As a result, correct reasoning could be scored as incorrect if the final answer is not presented in the expected format (e.g., missing/malformed \textbackslash boxed\{\} expressions). This creates a disconnect between reasoning quality and the actual evaluation performance of the model.
    \item \textbf{Strict Evaluation Criteria for Free-Form Answers:} Because free-form problems require exact matches for all sub-answers, there is no partial credit. What this means is that for minor numerical differences such as rounding, incorrect ordering of answers, or a single sub-answer missing from the final answer, the entire prediction is marked incorrect. The lack of leniency makes the task particularly challenging when it comes to multi-part problems.
    \item \textbf{Limited Adaptability (No Training):} Because the competition restricts the use of alternative models and discourages external tools, we cannot fine-tune the model or incorporate external computation (such as symbolic solvers). This limits the system's ability to correct systematic reasoning errors or adapt to specific domains within the dataset.
    \item \textbf{Truncation of Long Reasoning Chains:} Although large token limits are used, the model may still generate long reasoning traces that approach or exceed these limits. In cases like this, outputs might be truncated before a final answer is produced. This leads to evaluation failures despite partial reasoning process.
\end{itemize}

In order to address these limitations, we incorporate several practical fallback mechanisms within our inference pipeline.
\begin{itemize}
    \item \textbf{Adaptive Multi-Pass Inference:} Instead of allocating maximum computation to every problem, we use a staged approach: a fast initial pass handles easy problems efficiently and additional passes with higher token budgets are applied only to uncertain cases. This staged approach reduces overall runtime while still allowing deeper reasoning when needed \citet{snell2024scaling}.
    \item \textbf{Token Budget Scaling:} Different inference phases use different token limits. Lower limits are used for initial passes to reduce runtime, while higher limits are reserved for later stages to reduce the risk of truncation and allow more complete reasoning.
    \item \textbf{Structured Prompting:} We employ a structured reasoning prompt to guide the model toward more consistent and complete reasoning. This reduces common issues such as missing final answers or incomplete solution steps.
    \item \textbf{Output Formatting Constraints:} To ensure compatibility with the evaluation, prompts explicitly enforce formatting requirements such as placing the final answer in \textbackslash boxed\{\}, ordering multiple answers correctly, and separating reasoning from final output. This reduces evaluation errors caused by formatting inconsistencies.
    \item \textbf{Controlled Use of Self-Consistency:} For difficult problems, we optionally generate multiple independent outputs and select the most frequent answer. While this improves robustness, it is applied selectively due to its high computational cost \cite{wang2022selfconsistency}.
    \item \textbf{Uncertainty Detection via Output Heuristics:} We identify potentially unreliable outputs using simple heuristics, such as the absence of a clearly formatted final answer, unusually short or incomplete responses, or outputs that terminate due to token limits. These signals are used to trigger higher-budget inference phases within the adaptive pipeline. 
\end{itemize}

Overall, our approach mitigates key limitations - such as truncation, formatting sensitivity, and computational cost - through structured prompting and adaptive inference. While these strategies improve robustness under fixed-model constraints, challenges such as strict evaluation criteria and high runtime remain significant factors in overall system performance.

\subparagraph{Summary} Our approach differs from standard training-based methods by focusing entirely on inference-time optimization. This is necessary due to competition constraints that fix the underlying model. While this limits representational improvements, it allows us to leverage the strong reasoning capabilities of pretrained LLMs and improve performance through better prompting and decoding strategies.

% % % === Experiments
\section{Experiments}
\label{sec:experiments}
% 
% \textcolor{red}{In this section, clearly describe your empirical exploration of the problem and associated results. Provide evidence for your claims.
% Analyse and discuss your results.
% Note: Whenever you make claims such as \textit{Approach A performs ``better'' than approach B}, please be very specific in terms of \emph{what} it is better at and by \emph{how much}. Be transparent if it is only better on certain metrics, but not all.}

We conducted a series of controlled experiments to evaluate the impact of prompting and inference strategies on mathematical reasoning performance. Specifically, we study the effect of:
\begin{itemize}
    \item \textbf{Prompting Strategy:} Comparing naive prompting against structured prompting.
    \item \textbf{Inference Strategy:} Comparing single-pass generation with adaptive multi-phase inference, where additional computation is applied only to uncertain predictions.
    \item \textbf{Sampling Strategy:} Evaluating the effect of self-consistency (multiple generations with majority voting). Due to runtime constraints, this is only explored in small-scale environments.
\end{itemize}
All experiments were conducted on random subsets of the public dataset. Random sampling ensures coverage across different problem types while keeping runtime manageable.

\subsection{Baselines}
\label{sec:baselines}
% \textcolor{red}{Describe all the baselines, excluding the starter code. You should always start with simple models (such as a multi-layer Perceptron or linear regression) and gradually increase the complexity. You may use the starter code as a baseline, too.
% Use an itemized list to briefly summarize each of the models, their architecture, and parameters, and provide the correct reference if possible.}

We compare with the following baselines:
\begin{itemize}
    \item \textbf{Baseline A - Naive Prompting:} A single-pass generation using a minimal prompt without structured reasoning or formatting constraints. This baseline does not enforce structured reasoning or output formatting. It serves as a lower bound on performance.
    \item \textbf{Baseline B - Starter Configuration:} The provided starter pipeline using the required model (Qwen3-4B-Thinking-2507) with thinking-enabled generation and a large token budget. This baseline performs single-pass inference without structured prompting or adaptive compute allocation.
    \item \textbf{Our Method - Structured + Adaptive Inference:} Our approach augments the baseline with structured reasoning prompts, explicit output formatting constraints, adaptive multi-phase inference, and optional self-consistency.
    % Etc.
\end{itemize}

These baselines allow us to evaluate how each component contributes to the overall performance of the model.
% 
\subsection{Evaluation}
\label{sec:evaluation}
% \textcolor{red}{Describe how you evaluate your model and baselines. There is no need to provide long explanations for the Kaggle evaluation, but clearly explain any complementary evaluations you explore in your project (if any).}

We evaluate model performance using exact match accuracy on the public dataset.
\begin{itemize}
    \item \textbf{Multiple-Choice (MCQ):} A prediction is considered correct if the selected answer letter matches the ground truth.
    \item \textbf{Free-form:} A prediction is considered correct if and only if all required sub-answers exactly match the ground truth. No partial credit is awarded.
\end{itemize}

We report MCQ accuracy, Free-form accuracy, and overall accuracy. This strict evaluation reflects the competition scoring rules and emphasizes both reasoning correctness and output precision. The public dataset is used for validation and experimentation, while the private dataset is reserved for final submission and does not contain ground truth labels.

% 
\subsection{Implementation details}
\label{sec:implementation_details}
% \textcolor{red}{The goal of this section (together with other relevant parts of your paper) is to ensure full reproducibility without having to check your code.
% Some relevant points to include:
% \begin{itemize}
%     \item What computational platform/GPU did you use for training and testing? How long does it take to train your model for one epoch (going through the entire training data set once)?
%     \item Describe hyperparameters/regularization techniques/optimization details: What is your optimizer and learning rate, etc.? How many epochs? What is your batch size? Which HPs were tuned, and what range of values did you try? Was the tuning any different between baselines and your model? What were the final values that you used?
%     \item  How did you choose them? Provide any intuition of these choices: Explain why you made these design choices. Was it motivated by your past experience? Or was it due to the limitations of your computational platform? Be transparent if any hyperparameters are inconsistent between different runs (e.g., baseline A and B were trained with different learning rates).
% \end{itemize}
% }

\subparagraph{Model and Environment} We use the pretrained reasoning model Qwen3-4B-Thinking-2507. All experiments are conducted using GPU-accelerated inference via Hugging Face Transformers or vLLM. While vLLM is preferred for efficiency, Hugging Face Transformers is used as a fallback in the instance vLLM does not work.

\subparagraph{Generation Configuration} To support long reasoning chains, we use:
\begin{itemize}
    \item \textbf{Maximum Tokens:} 8192
    \item \textbf{Temperature:} 0.6
    \item \textbf{Top-p:} 0.95
    \item \textbf{Top-k:} 20
    \item \textbf{Repetition Penalty:} 1.0
\end{itemize}
These values follow the recommended decoding configuration from \citet{qwen3} and balance reasoning depth with output diversity.

\subparagraph{Inference Procedure} We perform sequential generation, which involves processing one problem at a time. This design not only simplifies debugging but also allows for progress tracking and preserves partial results if execution is interrupted. For our method, inference is performed using an adaptive multi-phase strategy where all questions receive an initial low-cost pass and uncertain predictions trigger additional higher-budget passes. As stated previously in section 2, uncertainty is detected using heuristics such as missing final answers, truncated outputs, or inconsistent reasoning patterns.
\subparagraph{Validation Setup} Due to high computational cost, we evaluate on a:
\begin{itemize}
    \item \textbf{10-question random subset} as an initial smoke test,
    \item \textbf{50-question subset} as ongoing validation.
\end{itemize}

Random sampling ensures coverage across problem types while keeping the runtime manageable.

\subparagraph{Runtime Considerations} Inference with a thinking model is computationally expensive, as it is observed to take $\sim$3-8 minutes per question. With that figure in mind, full dataset runs can take hours or days. As a result, full public evaluation is avoided during iteration and self-consistency is disabled during large runs.
% 
\subsection{Results}
\label{sec:results}
% \textcolor{red}{Present the results from starter code and preliminary implementation from your idea. This can be a mix of quantitative results (e.g., validation metrics or leaderboard score) and qualitative results (e.g., example visualizations of predictions, case studies, etc.). Use figures and tables. You are welcome to include any negative findings/results.}

\subparagraph{Initial Baseline Performance:} On a 10-question random subset of the public dataset:
% \begin{tab}[]
%     \centering
%     \begin{tabular}{c|c}
%         Metric & Accuracy \\
%         MCQ & 50.00\% ($\frac{2}{4}$)
%     \end{tabular}
%     \caption{Caption}
%     \label{tab:placeholder}
% \end{table}
\begin{center}
    \def\arraystretch{2.0}
    \begin{tabular}{||c c c c c||}
        \hline
        \textbf{Iteration} & \textbf{Method} & \textbf{MCQ} & \textbf{Free-form} & \textbf{Overall}\\
        \hline
         1 & Starter Baseline (N=10) & 50.0\% & 33.3\% & 40.0\% \\
        \hline
         2 & Structured Prompting (N=50) & 40.0\% & 36.7\% & 38.0\% \\
        \hline
         3 & Adaptive Multi-Phase Inference (N=50) & 55.0\% & 33.3\% & 42.0\% \\
        \hline
    \end{tabular}
\end{center}

These results demonstrate that naive prompting and single-pass inference are insufficient for reliable mathematical reasoning.
\subparagraph{Key Observations}
\begin{enumerate}
    \item \textbf{Structured Prompting Improves Free-Form Performance:} Comparing Iteration 1 and Iteration 2, free-form accuracy increased from 33.3\% to 36.7\%. This suggests that structured reasoning and improved formatting provide benefits for multi-answer problems.
    \item \textbf{Adaptive Inference Significantly Improves MCQ Performance:} Comparing Iteration 2 and Iteration 3, MCQ performance increased from 40.0\% to 55.0\%. This indicates that the additional reasoning budget helps refine discrete-choice predictions.
    \item \textbf{Overall Performance Improves with Adaptive Inference:} The overall accuracy increasing from 38.0\% in Iteration 2 to 42.0\% in Iteration 3 demonstrates that inference-time strategies can yield measurable improvements without modifying model weights.
    \item \textbf{Free-form Tasks Remain Challenging:} Free-form accuracy remains relatively flat, only fluctuating once from 33.0\% to 36.7\% before falling back down to the former value. This highlights persistent challenges, mainly related to multi-answer formatting, numerical precision, and strict exact-match evaluation.
\end{enumerate}

\subparagraph{Adaptive Inference Behavior} In the adaptive multi-phase setting using in Iteration 3, a substantial portion of problems required escalation to higher reasoning budgets. Specifically, only a small fraction of questions were resolved in the initial low-cost phase, while 52.0\% of the questions required full Phase 3 inference. This indicates that many problems demand deeper reasoning than provided by a single-pass approach. Despite this additional computation, 38.0\% of final outputs were still truncated. This suggests that even large token budgets are insufficient for some long reasoning chains. Furthermore, 46.0\% of questions remain uncertain after all phases. This highlights a limitation in the model's ability to consistently converge on confident answers. These results suggest that while adaptive inference improves overall performance, particularly for multiple-choice tasks, it is constrained by truncation effects and inherent model limitations for complex or multi-step problems.

\subparagraph{Error Analysis} To better understand model behavior, we analyze common failure modes:
\begin{enumerate}
    \item \textbf{Truncation Errors:} In several cases, the model produced long reasoning traces but failed to output a final answer due to token limits. These responses were marked incorrect despite containing partially correct reasoning.
    \item \textbf{Formatting Errors:} Even when the model arrived at the correct answer, missing or malformed \textbackslash boxed\{\} formatting causes incorrect evaluation. This highlights the importance of enforcing output structure.
    \item \textbf{Multi-Answer Failures:} For free-form problems with multiple required answers, the model often produced partially correct outputs. However, due to strict evaluation, any missing or incorrect sub-answer resulted in a zero score.
    \item \textbf{Numerical Precision Issues:} Small floating-point differences (such as rounding discrepancies) lead to mismatches with ground truth, particularly in problems involving continuous values.
\end{enumerate}


% \paragraph{Result 1.}
% \paragraph{Result 2.}
% etc.
% 
% \subsection{Ablations}
% \label{sec:ablations}
% Try to explore which design choices contribute to your final performance by removing/replacing individual components. Note that ablations are most convincing when you isolate the contribution of those individual components. For example, if you are comparing architecture A and architecture B, make sure to use the same hyperparameters as much as possible (or tune them for each architecture separately and compare the best version of each). Be clear but concise.
% 
% % % === Discussion
\section{Discussion}
\label{sec:discussion}
This project presented a pipeline for mathematical question answering, built around the Qwen/Qwen3-4B-Thinking-2507 LLM. The system addresses both free-form and multiple-choice mathematical reasoning problems from high-school through graduate-school level of difficulty. Since competition rules fix the model and prohibit external tools at inference time, our work focuses entirely on prompt engineering, adaptive compute allocation, and output formatting--techniques which we describe in detail in Section \ref{sec:methods}. At the milestone, we have a fully functional end-to-end pipeline (Contributions 1, 6, 7--as denoted in Section \ref{sec:intro}) and three scored evaluation iterations on small public subsets, with overall accuracy at 42\% on our most recent 50-question evaluation.

\subsection{Achievements.}
So far, we have achieved a complete, functional inference pipeline rather than a high-accuracy system. We are candid that the accuracy numbers reported are modest, and that they are based on a limited 50-question subset of the 1,126-question set due to severe computational constraints. Our contributions at this stage fall primarily in pipeline development and algorithm design, with accuracy improvements expected to follow as we complete full evaluations and iterate further. With that caveat, our preliminary results across three iterations are:
\begin{center}
    \def\arraystretch{2.0}
    \begin{tabular}{||c | c c c||}
        \hline
        \textbf{Iteration} & \textbf{MCQ} & \textbf{Free-form} & \textbf{Overall}\\
        \hline
        1 & 50.0\% ($\frac{2}{4}$) & 33.3\% ($\frac{2}{6}$) & 40.0\% ($\frac{4}{10}$)\\
        \hline
        2 & 40.0\% ($\frac{8}{20}$) & 36.7\% ($\frac{11}{30}$) & 38.0\% ($\frac{19}{50}$)\\
        \hline
        3 & 55.0\% ($\frac{11}{20}$) & 33.3\% ($\frac{10}{30}$) & 42.0\% ($\frac{21}{50}$)\\
        \hline
    \end{tabular}
\end{center}
The 42\% overall accuracy reflects the fundamental difficulty of mathematical reasoning for a 4B-parameter model operating without fine-tuning, external tools, or code execution--limitations that we identify in Section \ref{sec:methods}. The relative improvement from 38\% to 42\% between Iterations 2 and 3--particularly the 15 percentage point MCQ gain--is encouraging as a directional signal, but we acknowledge that the 50-question sample size limits our confidence in these trends. 

At this stage, our contributions are concentrated in pipeline development (Contributions 1, 5, 6, 7) and algorithm design (Contributions 2, 3, 4). We have not yet contributed to dataset curation or model development (no fine-tuning has been attempted). Full-set evaluation results, once the current run completes, will provide a more reliable accuracy picture and inform whether fine-tuning should be pursued.

\subsection{Method Description.}
The technical details of our approach are described in Section \ref{sec:methods}, so here we focus on design decisions and tradeoffs that emerged during implementation, as well as specifics of data handling and deviations from the starter code.

For starters, we found that Transformers sequential generation required approximately $\sim$3–8 minutes per question on our hardware, making any full-dataset evaluation infeasible (estimated 56–150 hours for the public set alone). Switching to vLLM on WSL2 reduced per-question latency to approximately 20–40 seconds. 

We initially used max$\_$tokens=2048 and observed systematic truncation, which are responses cut off before producing an extractable answer. We increased to max$\_$tokens=8192, which improved answer extractability but did not eliminate the problem. In our Iteration 3 evaluation, 38\% of final responses (19/50) were still truncated at the length limit, even after Phase 3's maximum budget. This was the single largest remaining source of answer loss.

We originally planned uniform $N=8$ self-consistency for all questions. At approximately 8 minutes per question, this would require over 150 hours for the public set. The adaptive multi-phase design (Contribution 3) was a direct response to this constraint. In practice, phase usage on our 50-question evaluation was:
\begin{center}
    \def\arraystretch{2.0}
    \begin{tabular}{||c | c c||}
        \hline
        \textbf{Phase} & \textbf{Count} & \textbf{Pct}\\
        \hline
        Phase 1 (resolved) & 7 & 14.0\% \\
        \hline
        Phase 2 (escalated) & 17 & 34.0\% \\
        \hline
        Phase 3 (max budget) & 26 & 52.0\% \\
        \hline
        Still uncertain after all phases & 23 & 46.0\% \\
        \hline
    \end{tabular}
\end{center}
The fact that only 14\% of questions resolved at Phase 1 and 46\% remained uncertain after all three phases revealed that the majority of problems in this dataset were genuinely difficult for a 4B model, and there is likely a ceiling on what inference-time methods alone can achieve without model improvement.

The public set (1,126 problems: 375 MCQ, 751 free-form) served as the validation set, and the private set (943 problems, no ground truth) is reserved for submission. We validated on random 50-question subsets drawn with RANDOM$\_$SEED=42 for reproducibility. No external data, data augmentation, or synthetic data was used. Input normalization was minimal, as problems were parsed as-is from JSONL with [ANS] placeholders counted programmatically. Output normalization mirrored the provided judger.py evaluator: LaTeX alias normalization, \textbackslash boxed\{\} pattern extraction, and MCQ letter extraction from traces.

Our implementation departs from the starter notebook in five respects, corresponding to Contributions 2–7: 
\begin{enumerate}
    \item Structured prompt templates replacing the generic baseline prompt,
    \item adaptive multi-phase generation replacing single-pass inference,
    \item checkpoint/resume for fault tolerance during multi-day runs,
    \item finish-reason diagnostics to detect and respond to truncation,
    \item and explicit multi-answer handling with [ANS]-counting logic.
\end{enumerate}
The starter code's basic workflow is preserved as the skeleton, with our contributions layered on top.

\subsection{Lessons Learned.}
Studying judger.py revealed it is forgiving on symbolic equivalence but strict on answer extraction format. This shifted our priority toward \textbackslash boxed\{\} compliance (Contribution 5) before investing in reasoning improvements. The progression from 38\% to 42\% showed that prompt design and compute allocation alone can improve accuracy (Contributions 2, 3). However, 46\% of questions remained uncertain after all three phases, suggesting a ceiling for inference-time methods on a 4B model without fine-tuning. We initially used max$\_$tokens=2048 and observed systematic truncation. Even after increasing to 8192 (Contribution 4), 38\% of responses were still truncated, and this is not just a quality issue but a correctness one, since missing \textbackslash boxed\{\} answers are scored as zero. MCQ accuracy improved significantly across iterations (50\% $\rightarrow$ 40\% $\rightarrow$ 55\%), while free-form remained flat ($\sim$33\%). Adaptive retries helped MCQ where reasoning is directionally correct, and free-form failures stem from formatting, precision, and multi-answer issues that our current approach does not fully address. Infrastructure work (Contribution 6) consumed significant early time but was essential, as without vLLM, the experiments would take days. Even with it, a full-dataset adaptive run exceeds 28 hours on our RTX 3070, fundamentally limiting iteration speed. If given more time at this stage, we would want to complete the full public-set evaluation for a statistically reliable baseline, conduct per-question error analysis to categorize failures by root cause (reasoning vs. formatting vs. truncation vs. tolerance), experiment with token budgets beyond 8192 to reduce truncation, and begin QLoRA \cite{hu2021lora}, \cite{dettmers2023qlora} fine-tuning on external math corpora such as OpenMathInstruct-2 \cite{toshniwal2024openmath2} and NuminaMath-CoT \cite{numina2024aimo}.

\subsection{Bottlenecks and Mitigation.}
Inference latency was the most critical bottleneck. Sequential Transformers generation required 3–8 minutes per question, making full evaluations infeasible. vLLM on WSL2 (Contribution 6) reduced latency by $\sim10\times$, but the full adaptive pipeline still took over 28 hours for 1,126 questions, limiting us greatly. GPU memory on the RTX 3070 (8 GB VRAM) constrained batch sizes and concurrent samples. We addressed this with conservative memory settings and conditional sample-count escalation (Contribution 3), which ultimately prevented OOM errors but limited throughput. Token truncation remains the largest source of answer loss. Despite increasing max$\_$tokens to 8192 and implementing truncation detection (Contribution 4), 38\% of responses in our latest run were still cut off before producing a final \textbackslash boxed\{\} answer. Strict free-form evaluation awards zero credit if any sub-answer is wrong, which partly explains the persistent MCQ–free-form accuracy gap despite formatting improvements.

\subsection{Plans Forward.}
Between now and the final report, we plan to:
\begin{itemize}
    \item Complete the full-dataset evaluation, which will serve as the true baseline for all subsequent work.
    \item Run full adaptive inference on all 943 private questions to produce the first leaderboard CSV, providing an external accuracy signal that complements our public validation results.
    \item Address truncation, as with 38\% of responses truncated and 46\% of questions uncertain after all phases, truncation is the single largest identified accuracy bottleneck. We plan to experiment with higher token budgets (16384+), prompt-level conciseness instructions to reduce reasoning verbosity, and tighter thinking$\_$budget constraints.
    \item Close the free-form accuracy gap, as the persistent gap between MCQ (55\%) and free-form (33.33\%) accuracy suggests that format-specific interventions are needed. We plan to develop specialized prompts for different free-form answer types (numerical, algebraic, interval, multi-part) and improve multi-answer formatting instructions.
    \item Categorize failures from full public results into reasoning errors, formatting failures, truncation-induced losses, and numeric tolerance mismatches to prioritize the highest-impact interventions for the final report.
    \item Potentially explore lightweight fine-tuning on external math reasoning datasets (OpenMathInstruct-2, NuminaMath-CoT).
\end{itemize}

\subsection{Team Responsibilities.}
\begin{center}
\def\arraystretch{1.5}
\begin{tabular}{||c | c | c||}
    \hline
    \textbf{Member} & \textbf{Background / Expertise} & \textbf{Role \& Contributions}\\
    \hline
    Tia Irani & CS major, experience with computer systems, optimization & Report writing \\
    \hline
    Leonel Alejandro Montoya & Math-CS major, experience with ML & Report writing \\
    \hline
    Diego Osborn & DS major, experience with statistical modeling/ML & Report writing \\
    \hline
    Sardor Sobirov & DS major, experience with full-stack dev/ML & Led pipeline development\\
    \hline
\end{tabular}
\end{center}

\subsection{Timelines.}
Estimate the timeline with the milestones to achieve your expected outcome in this class, considering the discussion provided above.
\begin{itemize}
    \item Week 2: Project setup
    \item Week 3: Literature survey and strategy report
    \item Week 4: Starter code reproduction and Iteration 1
    \item Week 5: Structured prompting (Iteration 2) and adaptive inference design
    \item Week 6: Adaptive inference evaluation (Iteration 3) and milestone report
    \item Week 7: Full-dataset results and error analysis
    \item Week 8: Truncation mitigation and free-form improvements
    \item Week 9: QLoRA fine-tuning exploration
    \item Week 10: Final evaluation, private submission, and report
\end{itemize}

\section{LLM Usage}
LLMs served as general-purpose assistants for research, implementation, and writing throughout this project. The team was responsible for all experimental design decisions, result interpretation, and validation of correctness. We used ChatGPT Deep Research to produce a literature survey on single-GPU math reasoning pipelines, covering model selection, prompting strategies, evaluator analysis, and inference-time scaling techniques. All strategic decisions (e.g., adopting adaptive multi-phase inference, selecting token budgets) were made by team members based on findings from this survey and our own experimentation. LLM-based coding assistants (Cursor/Claude) were used extensively for implementing the inference pipeline, including prompt construction, adaptive generation logic, checkpointing, answer extraction, and CSV export. In many cases, we described desired functionality in natural language and iteratively refined the generated code. All code was reviewed, tested, and validated by team members against the competition evaluator before use. LLMs assisted with drafting and editing portions of this report. All content was reviewed, revised, and verified by team members for accuracy and completeness. LLMs were consulted for environment setup (WSL2, vLLM, CUDA configuration), debugging generation parameters, and interpreting model card documentation for Qwen3-Thinking.

% % % === ACKNOWLEDGMENTS AND REFERENCES START
% \section*{Acknowledgements}
% This work was supported $\dots$.
\bibliography{references.bib}
\bibliographystyle{plainnat}  % Can also use abbrvnat, unsrt, etc.
% 
% % % === APPENDIX (optional)
% \newpage
% \appendix
% \section*{Appendix}
% \section{Appendix topic 1}
% 
\end{document}